\chapter{Log Levels and Intent}
Levels are filters, not decorations. They help people skim, and they help alerts route. A misused level creates noise or hides risk.

\section{Default Mapping}
Use a small set of levels. Keep meaning consistent across services.

\begin{table}[h]
\centering
\begin{tabularx}{\textwidth}{lX}
\toprule
Level & Use it when \\
\midrule
DEBUG & You need internal state for diagnosis. Do not ship to production unless controlled. \\
INFO & The event is expected and the outcome is normal. \\
WARN & The outcome is unusual but the request still completes. \\
ERROR & The event failed or data could not be processed. \\
FATAL & The process cannot continue safely. \\
\bottomrule
\end{tabularx}
\end{table}

\section{Outcome Drives Level}
The outcome should drive the level. If the outcome is \texttt{success}, the level is usually info. If it is \texttt{failure}, use error. If the system recovered, use warn. Rare exceptions are fine, but they should be explicit.

\section{Checklist: Level Hygiene}
\begin{itemize}
\item Every error-level event has an outcome of \texttt{failure}.
\item Warnings explain why recovery happened.
\item Debug logs are safe to disable without hiding important events.
\item Fatal logs are rare and always include a clear next step.
\end{itemize}

